\scanpage{304}
\mspage{45}
% archive sheet number: 303
est décrite de façon unique en disant qu'elle est \uline{constante}
\deleted{en dehors \deleted{de la} des intervalles \unsure{flèches|lignes} de transition}
sur chacune des ``composantes connexes'' (au sens de
la numérotation des sommets \unsure{\}|)} de $I_\tau$) de l'ens
\deleted{$\unsure{E}$}\inserted{$E'\deleted{\unsure{(f)}}=E'(f)\subset I_\tau=[0,n]$}
des sommets qui sont sources ou buts de flèches
de transition dégénérées pour $f$ i.e. de $\mathrm{Tr}(\tau)\setminus E$.
Et sur le $\alpha$.ième élém\textsuperscript{t} de $E$ \inserted{d'indice $i_\alpha$ dans $\mathrm{Tr}(\tau)$,}
\deleted{$(\unsure{1\leq i\leq r})$}
($\alpha\in[1,r=\operatorname{card} E]$), $f$ induit le foncteur
\diagram{\begin{tikzcd}[ampersand replacement=\&,column sep=small,row sep=small]
 \& J_{i_\alpha}\arrow[r]\arrow[d,dash] \& I_{\tau'}\\
 \unsure{\mathbb{\Delta}|\Delta}_{\varepsilon_\alpha}\arrow[ur,phantom,"\simeq" description,sloped] \&
 \substack{\text{intervalle}\\\text{de transition}\\\text{d'indice }i_\alpha\\\text{dans }I_\tau} \&
\end{tikzcd}}
donné par la \deleted{flèche} \uline{unique} \inserted{$\underline{\mathrm{Ty}}$-}\uline{flèche} de $I_\tau$
de $J_{i_\alpha}$ dans le $\alpha$.ième facteur, \inserted{de type} $\Delta_{\varepsilon_\alpha}^{n_\alpha}$, de $\tau'$\unsure{$*$|,}
donné par $(*)$
\[
 \Delta_{\varepsilon_\alpha}^{n_\alpha}\simeq\text{type de }\mathbb{\Delta}_{n_\alpha}^{\varepsilon_\alpha}\quad\ldots
\]

2\textsuperscript{o} \uline{Détermination de} $\operatorname{Ob}\,\underline{\mathrm{Ty}}_{/\tau'}$

Si $(\tau,f)\in\operatorname{Ob}\,\underline{\mathrm{Ty}}_{\unsure{\tau'|/\tau'}}$, \unsure{donc}
\[
 f\colon I_\tau\longrightarrow I_{\tau'},
\]
on lui associe \deleted{\unsure{donc}} d'abord la partie
\[
 \unsure{\mathrm{Cr}|\mathrm{Tr}}(\tau',f)=E'_f\inserted{{}=E'}\subset I_{\tau'}
\]
\inserted{formée} des sommets \inserted{(dits ``critiques pour $f$'', cf page 41)}
de $I_{\tau'}$ \deleted{\unsure{critiques} qui sont séparés,}
\inserted{\unsure{disons en l'un|disons en un} bout d'une flèche $f(u)$, où $u\in\mathrm{Tr}(\tau,f)$,}
% slanted
\deleted{c'est une partie de $I_{\tau'}$ contenant les sommets critiques}
i.e. \inserted{\unsure{$f$ croît strictement|$f$ croisse strictement}} \unsure{$f(u)$|$f u$} n'est pas dégénérée.
\deleted{\unsure{Celle|Cette} \unsure{Celle|Cette}} $E'$
contient \deleted{\unsure{en particulier les} $f$\unsure{en} $+$}
$f(\deleted{\unsure{E}} I_{\tau\,\mathrm{red}})$ (i.e.
les images par $f$ des sommets critiques de $I_\tau$)
